% !TeX root = ../se200_identity_regimes.tex

\section{Preliminaries}
\label{sec:preliminaries}

This paper refines the referential-regime component of the neutral-substrate
constraint introduced in SE-100. In SE-100, a referential regime consists of
individuation, co-reference, and persistence conditions. The present paper
analyzes the persistence and transformation behavior needed for those
referential commitments to remain usable as common ground under persistent
interpretive disagreement.

This section fixes the formal vocabulary used in the derivations.
Identity is represented via regime-specific equivalence relations
and the quotient structures induced by those relations.
Statements are formulated in terms of equivalence classes;
the derivations in this section rely on invariance under admissible transformations.



%--------------------------------------------
\subsection{Neutral Substrate Constraint}
\label{subsec:ns_constraint}

\begin{definition}[Neutral substrate]
  \label{def:neutral_substrate}
  A neutral substrate is a representation layer
  that does not encode downstream interpretive content.
\end{definition}

\begin{definition}[Extension]
  \label{def:extension}
  An extension is any admissible addition, refinement, or reorganization
  of representations that preserves neutrality.
\end{definition}

\begin{definition}[Extension stability (extension invariance)]
  \label{def:extension_stability}
  A condition is extension-stable (equivalently, extension-invariant)
  if it is preserved under all admissible extensions.
\end{definition}

\begin{definition}[Admissible ontology]
  \label{def:admissible_ontology}
  An ontology is admissible if it has explicitly specified
  identity and persistence conditions that are extension-stable.
\end{definition}


The neutral substrate represents entities without fixing downstream interpretation.
Identity and persistence must be stable under extensions
by heterogeneous and potentially conflicting interpretations.

%--------------------------------------------
\subsection{Representation Space}
\label{subsec:representation_space}

This subsection fixes the representation domain
on which all regimes and transformations are defined.

\begin{definition}[Representation space]
  \label{def:representation_space}
  Let $\mathcal{R}$ denote the set of all well-formed representations
  that satisfy neutrality and stable reference.
\end{definition}

%--------------------------------------------
\subsection{Transformation Basis}
\label{subsec:transformation_basis}

This subsection fixes the transformation families
that parameterize identity and persistence behavior.

\begin{definition}[Transformation family]
  \label{def:transform_family}
  Let $\mathbb{F}$ be a fixed set of transformation families:
  \[
    \mathbb{F}
    =
    \{ \mathrm{RE}, \mathrm{AN}, \mathrm{RF},
    \mathrm{AD}, \mathrm{RC}, \mathrm{RA},
    \mathrm{SU}, \mathrm{BF}, \mathrm{PV}, \mathrm{SE} \}.
  \]
  Each $f \in \mathbb{F}$ denotes a transformation family,
  i.e., a class of admissible extension-relative transformations on $\mathcal{R}$.
\end{definition}


%--------------------------------------------
\subsection{Regime Profiles}
\label{subsec:regime_profiles}

This subsection defines regime profiles as the structures that
determine identity, applicability, classification, and persistence.

\begin{definition}[Applicability]
  \label{def:applicability}
  For a regime profile $P$, the applicability map
  \[
    \mathrm{Appl}_P : \mathbb{F} \to \{0,1\}
  \]
  is defined by
  \[
    \mathrm{Appl}_P(f) = 1
  \]
  iff the transformation family $f$ is applicable under $P$.
\end{definition}

\begin{definition}[Transformation classification]
  \label{def:transformation_classification}
  For a regime profile $P$, the classification map is the function
  \[
    \mathrm{Class}_P :
    \{ f \in \mathbb{F} \mid \mathrm{Appl}_P(f)=1 \}
    \to
    \{ \mathrm{IGN}, \mathrm{PRS}, \mathrm{BRK} \}.
  \]
  Thus classification is defined only for applicable transformation families.
\end{definition}

For transformation families $f$ with $\mathrm{Appl}_P(f) = 0$,
classification is undefined.
In profile tables, such entries are marked N/A,
indicating inapplicability rather than an
unclassified applicable transformation.

\begin{definition}[Persistence condition]
  \label{def:persistence_condition}
  For a regime profile $P$, the persistence condition
  \[
    \mathsf{Persist}_P \subseteq
    \{ f \in \mathbb{F} \mid \mathrm{Appl}_P(f)=1 \text{ and } \mathrm{Class}_P(f)=\mathrm{PRS} \}
  \]
  specifies which applicable identity-preserving transformation families
  count as continuation of the same entity.
\end{definition}


  Identity is given by equivalence classes;
  persistence governs continuity under
  transformations classified as preserving identity.


\begin{definition}[Regime profile]
  \label{def:regime_profile}
  A regime profile is a tuple:
  \[
    P = (C_P, B_P, \mathrm{Appl}_P, \mathrm{Class}_P, \mathsf{Persist}_P)
  \]
  where:
  \begin{itemize}
    \item $C_P$ is the carrier of identity,
    \item $B_P$ is the identity basis,
    \item $\mathrm{Appl}_P : \mathbb{F} \to \{0,1\}$ is the applicability map,
    \item $\mathrm{Class}_P$ is a function assigning each applicable transformation family
          a value in $\{ \mathrm{IGN}, \mathrm{PRS}, \mathrm{BRK} \}$,
    \item $\mathsf{Persist}_P$ specifies persistence behavior for transformations
          classified as preserving identity.
  \end{itemize}
\end{definition}

\begin{definition}[Regime]
  \label{def:regime}
  A regime $R$ implements a regime profile $P$ if the
  identity basis, applicability map, transformation classification,
  and persistence behavior of $R$ are exactly those specified by $P$,
  i.e., $R$ realizes $B_P$, $\mathrm{Appl}_P$,
  $\mathrm{Class}_P$, and $\mathsf{Persist}_P$ on $\mathcal{R}$.
\end{definition}



%--------------------------------------------
\subsection{Admissible Transformations}
\label{subsec:admissible_transformations}

This subsection defines admissibility conditions
for transformations relative to a regime profile.

\begin{definition}[Admissible transformation]
  \label{def:admissible_transformation}
  A transformation $\tau : \mathcal{R} \to \mathcal{R}$ is admissible under a regime profile $P$ if:
  \begin{enumerate}
    \item $\tau$ preserves well-formedness,
    \item $\tau$ preserves neutrality,
    \item $\tau$ preserves the invariants determined by the identity basis $B_P$,
      where $B_P$ is the identity basis component of the regime profile $P$
      defined in Definition~\ref{def:regime_profile}.
  \end{enumerate}
\end{definition}

\begin{definition}[Admissible transformation class]
  \label{def:admissible_transformation_class}
  For a regime profile $P$, let $T_P$ denote the set of all transformations
  admissible under $P$.
\end{definition}


  No closure or invertibility assumptions are imposed on $T_P$
  unless explicitly stated.


%--------------------------------------------
\subsection{Identity and Quotients}
\label{subsec:identity_quotients}

This subsection defines identity as an equivalence relation induced by
transformation classification and introduces the associated quotient structure.

\begin{definition}[Induced identity]
  \label{def:induced_identity}
  Let $P$ be a regime profile.
  Define $\sim_P$ to be the equivalence relation on $\mathcal{R}$
  generated by the transformation families $f \in \mathbb{F}$ such that
  \[
    \mathrm{Appl}_P(f)=1
    \quad\text{and}\quad
    \mathrm{Class}_P(f)=\mathrm{PRS}.
  \]
\end{definition}

\begin{definition}[Equivalence class]
  \label{def:equivalence_class}
  The identity of $x$ under $P$ is its equivalence class:
  \[
    [x]_P = \{ y \in \mathcal{R} \mid x \sim_P y \}.
  \]
\end{definition}

\begin{definition}[Quotient structure]
  \label{def:quotient_structure}
  The set of entities under profile $P$ is the quotient:
  \[
    \mathcal{E}_P := \mathcal{R} / \sim_P.
  \]
\end{definition}


  All identity statements are interpreted at the level of equivalence classes in $\mathcal{E}_P$.


%--------------------------------------------
\subsection{Invariance and Persistence}
\label{subsec:invariance_persistence}

This subsection characterizes invariance conditions
required for well-defined identity under a regime profile.

\begin{definition}[Invariant]
  \label{def:invariant}
  An invariant for a regime profile $P$ is a property preserved under all
  applicable transformation families classified as
  PRS by $\mathrm{Class}_P$.
\end{definition}


  Invariance under admissible transformations is a necessary condition
  for well-defined identity under a regime profile.



%--------------------------------------------
\subsection{Profile Comparison}
\label{subsec:profile_comparison}

This subsection defines equivalence and non-collapse
between regime profiles via their induced identity relations.

\begin{definition}[Profile equivalence]
  \label{def:profile_equivalence}
  Two regime profiles $P, Q$ are equivalent if they induce the same
  equivalence relation on $\mathcal{R}$, i.e.,
  \[
    \sim_P = \sim_Q.
  \]
  Equivalently, $P$ and $Q$ are equivalent if every identity class under $P$
  is also an identity class under $Q$, and vice versa.
\end{definition}


  If $P$ and $Q$ are equivalent, then
  $\mathcal{E}_P = \mathcal{E}_Q$.


\begin{definition}[Collapse]
  \label{def:collapse}
  Two regime profiles $P, Q$ collapse if they are equivalent but not identical,
  i.e., $\sim_P = \sim_Q$ but $P \neq Q$.
\end{definition}

\begin{definition}[Non-collapse]
  \label{def:non_collapse}
  Two regime profiles $P, Q$ are non-collapsing if they are not equivalent,
  i.e., $\sim_P \neq \sim_Q$.
  Equivalently, $P$ and $Q$ are non-collapsing if there exists at least one pair of representations
  that are in the same identity class under one profile but not under the other.
\end{definition}


%--------------------------------------------
\subsection{Transformation-Driven Structure}
\label{subsec:transformation_structure}

This subsection formalizes transformation patterns and the conditions
under which they are realized or force refinement of regime profiles.

\begin{definition}[Coherent transformation pattern]
  \label{def:coherent_transformation_pattern}
  A transformation pattern $\Pi$ is a pair
  \[
    \Pi = (\mathrm{Appl}_\Pi,\mathrm{Class}_\Pi)
  \]
  where
  \[
    \mathrm{Appl}_\Pi : \mathbb{F} \to \{0,1\}
  \]
  and
  \[
    \mathrm{Class}_\Pi :
    \{ f \in \mathbb{F} \mid \mathrm{Appl}_\Pi(f)=1 \}
    \to
    \{ \mathrm{IGN}, \mathrm{PRS}, \mathrm{BRK} \}.
  \]
  The pattern $\Pi$ is coherent if:
  \begin{itemize}
    \item it satisfies the neutrality constraint,
    \item it induces a well-defined equivalence relation on $\mathcal{R}$,
    \item it contains no internal contradictions in applicability or classification.
  \end{itemize}
\end{definition}

\begin{definition}[Profile realization]
  \label{def:profile_realization}
  A regime profile $P$ realizes a transformation pattern $\Pi = (\mathrm{Appl}_\Pi,\mathrm{Class}_\Pi)$ iff
  \[
    \mathrm{Appl}_\Pi = \mathrm{Appl}_P
    \quad\text{and}\quad
    \mathrm{Class}_\Pi = \mathrm{Class}_P.
  \]
\end{definition}

\begin{definition}[Split pressure]
  \label{def:split_pressure}
  A provisional regime profile $P$ is under split pressure if there exists a transformation
  family $f \in \mathbb{F}$ such that no single pair
  \[
    (\mathrm{Appl}_P(f), \mathrm{Class}_P(f))
  \]
  is consistent with all admissible cases assigned to $P$.
\end{definition}

\begin{definition}[Forced regime profile]
  \label{def:forced_regime_profile}
  A regime profile is forced if there exists a coherent transformation pattern $\Pi$
  required for stable identity under extension such that no existing profile realizes $\Pi$.
\end{definition}


  Since $\mathrm{Class}_P$ is a function on applicable
  transformation families, each applicable $f \in \mathbb{F}$ receives
  exactly one classification under $P$.
  Accordingly, no profile assigns more than one classification
  to the same transformation family.


%--------------------------------------------
\subsection{Minimal Regime Sets}
\label{subsec:minimal_sets}

This subsection defines minimal sets of regime profiles
sufficient to realize all required transformation patterns without collapse.

\begin{definition}[Minimal regime set]
  \label{def:minimal_regime_set}
  A set of regime profiles $\mathcal{P}$ is minimal if:
  \begin{itemize}
    \item every coherent transformation pattern required for stable identity under extension
          is realized by some $P \in \mathcal{P}$,
    \item no two profiles in $\mathcal{P}$ collapse,
    \item removing any profile from $\mathcal{P}$ violates one of the above conditions.
  \end{itemize}
\end{definition}

%--------------------------------------------
\subsection{Obstruction and Separation}
\label{subsec:obstruction}

This subsection introduces obstruction and witness concepts
used to establish non-collapse and necessity results.

\begin{definition}[Obstruction]
  \label{def:obstruction}
  An obstruction is a structural condition that prevents a target property
  from being satisfied under the admissibility constraints.
\end{definition}

\begin{definition}[Witness]
  \label{def:witness}
  A witness is a concrete configuration demonstrating the presence of an obstruction.
\end{definition}

\begin{definition}[Irreducible obstruction]
  \label{def:irreducible_obstruction}
  An obstruction is irreducible if no admissible transformation eliminates it.
\end{definition}

\begin{definition}[Separation witness]
  \label{def:separation_witness}
  A separation witness for regime profiles $P, Q$ is a pair $(x,y) \in \mathcal{R}^2$
  such that $x \sim_P y$ and $x \not\sim_Q y$, or vice versa.
\end{definition}

%--------------------------------------------
\subsection{Logical Form of Results}
\label{subsec:logical_form}

This subsection specifies the form
in which results are stated and proved.


  Statements of results are formulated
  in terms of quotient structures $\mathcal{E}_P$.
  Proofs proceed by establishing invariance properties
  under admissible transformation classifications.




%--------------------------------------------
%--------------------------------------------
\subsection{Representation Constraints}
\label{subsec:representation_constraints}

\subsubsection{Admissibility Constraints}

\begin{proposition}[Admissibility Boundary]
\label{prop:admissibility_boundary}
Systems requiring primitive causal, normative, or evaluative relations
are not representable under the admissible relation set.
\end{proposition}

\begin{proof}
Primitive causal and normative relations are excluded by admissibility.
\end{proof}

\begin{proposition}[Causal Non-Representability]
\label{prop:causal_nonrepresentability}
The representation preserves temporal order but does not represent causal production.
\end{proposition}

\begin{proof}
Only the admissible relation $\mathrm{precedes}$ is available,
which encodes order but not causal production.
\end{proof}

\begin{proposition}[Normative Non-Representability]
\label{prop:normative_nonrepresentability}
Normative force is not represented; only structural placement of normative entities is preserved.
\end{proposition}

\begin{proof}
No admissible relation encodes obligation, permission, or prohibition.
\end{proof}

\subsubsection{Representation Properties}

\begin{proposition}[REC Dependence]
\label{prop:rec_dependence}
Observability is mediated by REC-anchored paths.
\end{proposition}

\begin{proof}
REC nodes provide the only linkage to descriptive records within the representation.
\end{proof}

\begin{proposition}[Granularity Dependence]
\label{prop:granularity_dependence}
Path structure depends on modeling granularity.
\end{proposition}

\begin{proof}
Refinement alters vertex and edge structure, thereby altering path families.
\end{proof}

\begin{corollary}[Path Sufficiency]
\label{cor:path_sufficiency}
All admissible structure is representable as path families over admissible relations.
\end{corollary}

\begin{proof}
Follows immediately from expressive adequacy of the path grammar;
% see Proposition~\ref{prop:path_grammar_adequacy} in
Section~\ref{sec:formal_setting}.
\end{proof}
